% Extended canonicalization rationale (moved out of main §\ref{sec:canonicalization}).

\section{Canonicalization policy: full rationale}
\label{app:canonicalization}

This appendix records the alternative SMILES-canonicalization policies we
considered, the empirical evidence against the merging policies, and the
chemical interpretation of that evidence.  The main text adopts \emph{Option
D} (plain RDKit canonical with all stereochemistry preserved); the
material below explains why.

\subsection{Candidate policies}

\begin{table}[h]
\centering\small
\caption{Candidate canonicalization policies for solute \SMILES{} and the
resulting unique-solute count when applied to the \numval{1\,525} raw
solute \SMILES{} in \textsc{BigSolDB}~v2.1.}
\label{tab:canonicalization}
\begin{tabular}{@{}clr@{}}
\toprule
Option & Definition & Unique solutes \\
\midrule
A & tautomer-enumerate + strip \texttt{@/@@} + strip \texttt{/\textbackslash} & 1\,506 \\
B & no tautomer         + strip \texttt{@/@@} + strip \texttt{/\textbackslash} & 1\,508 \\
C & no tautomer         + strip \texttt{@/@@}, keep \texttt{/\textbackslash} & 1\,510 \\
\textbf{D} & plain canonical, \texttt{isomericSmiles=True} (keep all stereo) & \textbf{1\,525} \\
\bottomrule
\end{tabular}
\end{table}

Options A--C all merge distinct compounds.  Two concrete failures of stereo
stripping illustrate the cost: \emph{fumaric} acid (trans,
\verb|O=C(O)/C=C/C(=O)O|) and \emph{maleic} acid (cis,
\verb|O=C(O)/C=C\C(=O)O|), which differ in aqueous solubility by roughly
$100\times$, become identical after removing
\texttt{/}/\texttt{\textbackslash}.  A concrete failure of tautomer
enumeration: a bicyclic norbornene anhydride is silently \emph{aromatized}
by the RDKit tautomer enumerator into a structurally nonsensical
``aromatic diol''.

\subsection{Empirical audit of Option-C merges}

To choose between the remaining options we audit every Option-C merge
group: for each group we find all $(\text{solvent}, T)$ cells in which
$\geq 2$ raw-SMILES members report measurements and compute
$|\Delta \log_{10} S|$ at matched conditions.  If enantiomers genuinely
had identical solubility in achiral solvents, these deltas should lie at
or below the aleatoric floor ($\sim 0.1$~log~S; §\ref{sec:aleatoric}).
We find the opposite: \emph{10 of 15} Option-C merge groups disagree by a
\emph{median} of \numval{0.29}--\textbf{\numval{1.31}}~log~S across
3--117 overlap pairs (Table~\ref{tab:stereo_audit}), i.e.\ 5--20$\times$
the aleatoric floor.

\begin{table}[h]
  \centering\small
  \caption{Empirical audit of Option-C stereo-merge groups.
    Median / max $|\Delta \log_{10} S|$ across all $(\mathrm{solvent}, T)$
    cells where $\ge 2$ raw-SMILES members of the merged group have
    measurements.  The four ``silent'' merges (bottom) have no overlap
    cells and cannot be empirically tested.}
  \label{tab:stereo_audit}
\begin{tabular}{@{}lrrrr@{}}
\toprule
Merge group (common name)              & total rows & overlap pairs & median $|\Delta|$ & max $|\Delta|$ \\
\midrule
Brassinolide (2 stereoisomers)         & 146 &  36 & \textbf{1.31} & 2.01 \\
D / L-psicose                          & 179 &  36 & 0.61 & 0.89 \\
D / L-malic acid                       & 141 &  21 & 0.52 & 0.99 \\
D / L-tryptophan                       & 229 &  84 & 0.43 & 0.97 \\
D / L-norbornene anhydride             & 245 & 117 & 0.40 & 0.56 \\
L-leucine / mixed                      & 126 &   8 & 0.39 & 0.45 \\
D / L-tyrosine                         & 123 &  16 & 0.37 & 1.07 \\
N-acetyl-methionine D / L              & 188 &  45 & 0.37 & 0.73 \\
Ibuprofen (S vs mixed)                 & 106 &   1 & 0.29 & 0.29 \\
\midrule
L-isoleucine (silent)                  & 126 &   0 &   --- &   --- \\
Ofloxacin (silent)                     & 124 &   0 &   --- &   --- \\
Naproxen (silent)                      & 115 &   0 &   --- &   --- \\
Tartaric acid (silent)                 &  49 &   0 &   --- &   --- \\
\bottomrule
\end{tabular}
\end{table}

\subsection{Chemical interpretation}

The main chemically interpretable explanation for the observed
disagreement is that enantiomeric compounds can be crystallized in
different polymorphic modifications, for which solubility differences
are commonly observed.  Other potential causes are
lab~$\times$~chirality-label correlation in measurement protocol, or
silent mislabelling of racemate as L (or vice versa) at the
\textsc{BigSolDB} extraction stage.  We cannot disambiguate these from
the data, but \emph{any} of the three implies that merging averages
away a $\sim 0.4$-log-unit signal we cannot recover.  Adopting Option~D
(stereo-preserving canonical \SMILES) keeps that signal in the labels;
chirality-blind featurizations then pay the cost in model error rather
than hiding it inside the label.
